\begin{lemma}
\label{lemma-normalization-same-after-completion}
Let $A$ be a reduced Nagata local ring of dimension $1$.
Let $A \to A'$ be as in Lemma \ref{lemma-pre-delta-invariant}.
Let $A^h$, $A^{sh}$, resp.\ $A^\wedge$
be the henselization, strict henselization, resp.\ completion of $A$.
Then $A^h$, $A^{sh}$, resp. $A^\wedge$ is a reduced Nagata local
ring of dimension $1$ and
$A' \otimes_A A^h$, $A' \otimes_A A^{sh}$, resp. $A' \otimes_A A^\wedge$
is the integral closure of $A^h$, $A^{sh}$, resp.\ $A^\wedge$
in its total ring of fractions.
\end{lemma}

\begin{proof}
Observe that $A^\wedge$ is reduced, see
More on Algebra, Lemma \ref{more-algebra-lemma-completion-reduced}.
The rings $A^h$ and $A^{sh}$ are reduced by
More on Algebra, Lemma \ref{more-algebra-lemma-henselization-reduced}.
The dimensions of $A$, $A^h$, $A^{sh}$, and $A^\wedge$ are the same
by More on Algebra, Lemmas
\ref{more-algebra-lemma-completion-dimension} and
\ref{more-algebra-lemma-henselization-dimension}.

\medskip\noindent
Recall that a Noetherian local ring is Nagata if and only if the
formal fibres of $A$ are geometrically reduced, see
More on Algebra, Lemma \ref{more-algebra-lemma-Nagata-local-ring}.
This property is inherited by $A^h$ and $A^{sh}$, see
the material in More on Algebra, Section
\ref{more-algebra-section-properties-formal-fibres}
and especially Lemma \ref{more-algebra-lemma-henselization-P-ring}.
The completion is Nagata by
Algebra, Lemma \ref{algebra-lemma-Noetherian-complete-local-Nagata}.

\medskip\noindent
Now we come to the statement on integral closures. Before continuing
let us pick $f \in \mathfrak m$ as in
Lemma \ref{lemma-pre-pre-delta-invariant}.
Then the image of $f$ in $A^h$, $A^{sh}$, and $A^\wedge$
clearly is an element satisfying properties (1) -- (6)
in that ring.

\medskip\noindent
Since $A \to A'$ is finite we see that
$A' \otimes_A A^h$ and $A' \otimes_A A^{sh}$
is the product of henselian local rings finite over $A^h$ and
$A^{sh}$, see Algebra, Lemma \ref{algebra-lemma-finite-over-henselian}.
Each of these local rings is the henselization of $A'$ at a
maximal ideal $\mathfrak m' \subset A'$ lying over $\mathfrak m$, see
Algebra, Lemma \ref{algebra-lemma-quasi-finite-henselization} or
\ref{algebra-lemma-quasi-finite-strict-henselization}.
Hence these local rings are normal domains by
More on Algebra, Lemma \ref{more-algebra-lemma-henselization-normal}.
It follows that $A' \otimes_A A^h$ and $A' \otimes_A A^{sh}$
are normal rings. Since $A^h \to A' \otimes_A A^h$ and
$A^{sh} \to A' \otimes_A A^{sh}$ are finite (hence integral) and since
$A' \otimes_A A^h \subset (A^h)_f = Q(A^h)$ and
$A' \otimes_A A^{sh} \subset (A^{sh})_f = Q(A^{sh})$
we conclude that $A' \otimes_A A^h$ and $A' \otimes_A A^{sh}$ are
the desired integral closures.

\medskip\noindent
For the completion we argue in entirely the same manner. First,
by Algebra, Lemma \ref{algebra-lemma-completion-finite-extension} we have
$$
A' \otimes_A A^\wedge = (A')^\wedge =
\prod\nolimits (A'_{\mathfrak m'})^\wedge
$$
The local rings $A'_{\mathfrak m'}$ are normal and have dimension $1$
(by Algebra, Lemma \ref{algebra-lemma-finite-in-codim-1} for example or
the discussion in
Algebra, Section \ref{algebra-section-homomorphism-dimension}).
Thus $A'_{\mathfrak m'}$ is a discrete valuation ring, see
Algebra, Lemma \ref{algebra-lemma-characterize-dvr}.
Hence $(A'_{\mathfrak m'})^\wedge$ is a discrete valuation ring
by More on Algebra, Lemma \ref{more-algebra-lemma-completion-dvr}.
It follows that $A' \otimes_A A^\wedge$ is a normal ring and
we can conclude in exactly the same manner as before.
\end{proof}